\section{Заключение}

В данной работе было проведено исследование одной из фундаментальных проблем прогнозирования
 хаотических систем — дефицита данных для надежной реконструкции аттрактора. Была экспериментально
  проверена гипотеза о возможности восполнения этого дефицита за счет привлечения данных из схожих 
  по динамике, но формально независимых систем.
Ключевым результатом работы стало выявление нетривиального характера 
прогностической совместимости временных рядов. Эксперименты показали, 
что положительный эффект от аугментации не гарантирован и зависит от
 сложной нелинейной согласованности динамик, а не от формальной близости 
 управляющих параметров. Пример с рядами при 
r=28.0001 и 
r=28.01 наглядно продемонстрировал несостоятельность стандартных метрик
 глобальной схожести и поставил задачу разработки более глубокого критерия отбора рядов-доноров.
Основным вкладом данной работы является предложение и апробация такого 
критерия, основанного на двухэтапном сравнении конечномерных аппроксимаций оператора 
Перрона-Фробениуса и расчете расстояния Шамфера между двумя рядами. Ряды, отобранные с его помощью, показали 
статистически существенно значимо лучшие результаты по сравнению со случайно выбранными, что
 подтверждено критерием Манна-Уитни. По крайней мере половина этих рядов была выгоднее,чем  выборка,состоящая из оригинального ряда
  по RMSE. Эксперименты на данных системы
  Лоренца показали, что действительно полезные ряды-доноры встречаются редко, что подчеркивает практическую
   ценность предложенного фильтрующего механизма. Более того, точность самих метрик 
   зависит от длины анализируемых рядов, что является естественным ограничением. 
   Тем не менее, предложенный подход превращает метод межсерийного расширения данных
    в управляемый и надежный инструмент для повышения полноты и точности прогноза в
     условиях, когда получение длинных стационарных выборок невозможно,справляясь с задачей трансферного обучения.
    Это создает перспективы для его применения в таких прикладных областях, как
     финансовые рынки, климатология,
 анализ биомедицинских сигналов или изучение паттернов энергопотребления \cite{PerezChacon2018}.


 \section{Благодарности}
Проведение вычислительных экспериментов, описанных в данной работе, стало
 возможным благодаря использованию ресурсов Суперкомпьютерного комплекса НИУ ВШЭ.